% !TeX TS-program = xelatex
\documentclass[11pt, a4paper]{article}

% ── Geometry ──
\usepackage[
  top=0.46in,
  bottom=0.38in,
  left=0.58in,
  right=0.58in
]{geometry}

% ── Chinese support ──
\usepackage{xeCJK}
\setCJKmainfont{FandolSong}[
  BoldFont=FandolHei-Bold,
  ItalicFont=FandolKai-Regular
]
\setCJKsansfont{FandolHei}[BoldFont=FandolHei-Bold]
\setCJKmonofont{FandolFang-Regular}

% ── Fonts (Latin) ──
\usepackage{fontspec}
% Use Computer Modern (default) for Latin text; override if you prefer:
% \setmainfont{TeX Gyre Termes}

% ── Packages ──
\usepackage{enumitem}
\usepackage{titlesec}
\usepackage{hyperref}
\usepackage{xcolor}

% ── Colors ──
\definecolor{linkblue}{HTML}{0563C1}
\definecolor{sectioncolor}{HTML}{000000}
\definecolor{rulecolor}{HTML}{333333}
\definecolor{entryrule}{HTML}{D8D8D8}
\definecolor{muted}{HTML}{444444}

% ── Hyperlinks ──
\hypersetup{
  colorlinks=true,
  urlcolor=linkblue,
  linkcolor=linkblue,
  pdfauthor={韩天硕},
  pdftitle={韩天硕 - 简历},
}

% ── No page numbers ──
\pagestyle{empty}

% ── Section style ──
\titleformat{\section}
  {\normalsize\bfseries\color{sectioncolor}}
  {}{0em}{}
  [\vspace{-0.6em}\textcolor{rulecolor}{\rule{\textwidth}{0.6pt}}]
\titlespacing{\section}{0pt}{6pt}{3pt}

% ── List style ──
\setlist[itemize]{
  leftmargin=1.25em,
  itemsep=0pt,
  parsep=0pt,
  topsep=0pt,
}

% ── Helper commands ──

% Resume entry: \resumeentry{Title}{Date}
\newcommand{\resumeentry}[2]{%
  \noindent\textbf{#1}\hfill #2\par\vspace{1pt}
}

% Resume sub-entry: \resumesub{Subtitle}
\newcommand{\resumesub}[1]{%
  \noindent{\color{muted}#1}\par\vspace{1pt}
}

% Publication entry: \pub{Title}{Authors}{Venue}
\newcommand{\pub}[3]{%
  \noindent\textbf{#1}\par
  \noindent{\small #2}\par
  \noindent{\small\color{muted}#3}\par
}

% Lightweight divider between dense resume entries.
\newcommand{\entrydivider}{%
  \par\vspace{1pt}
  \noindent\textcolor{entryrule}{\rule{\textwidth}{0.3pt}}\par
  \vspace{1pt}
}

% ── Reduce spacing ──
\setlength{\parindent}{0pt}
\setlength{\parskip}{0pt}
\setlength{\emergencystretch}{2em}

% ══════════════════════════════════════════
\begin{document}

% ── Header ──
\begin{center}
  {\LARGE\bfseries 韩天硕}\\[3pt]
  {\small
  (+86)\,173-3953-9195 \enspace$\mid$\enspace
  \href{mailto:hantianshuo@iie.ac.cn}{hantianshuo@iie.ac.cn} \enspace$\mid$\enspace
  \href{https://github.com/keymaker-arch}{github.com/keymaker-arch}}\\[4pt]
  {\normalsize\bfseries 安全研究员 / 漏洞挖掘与利用 / Linux内核安全}
\end{center}

\vspace{-7pt}

% ── 教育经历 ──
\section{教育经历}

\resumeentry{中国科学院大学，信息工程研究所\enspace|\enspace 网络空间安全，直博在读}{2021.09 -- 2027.01（预计）}
\begin{itemize}
  \item 研究方向：Linux内核安全、漏洞挖掘与利用、模糊测试、AI for Security。
  \item 代表成果发表于IEEE S\&P、USENIX Security、ACM CCS等世界顶级安全会议。
\end{itemize}

\resumeentry{中国农业大学，生命科学学院\enspace|\enspace 生物科学，理学学士}{2017.09 -- 2021.06}
\begin{itemize}
  \item 跨专业保送至中国科学院大学网络空间安全专业直博；曾获全国大学生数学建模竞赛国家级一等奖、CUPT青年物理学家竞赛华北赛区一等奖、iGEM国际基因工程机器大赛金奖。
\end{itemize}

% ── 学术成果 ──
\section{学术成果}

\pub{NetPanic: The Attack Surface You Can't Syscall}
{\textbf{Tianshuo Han}, Zong Cao, Zhen Dong, Xiapu Luo, Zhenyu Song, and Jian Liu}
{47th IEEE S\&P, 2026, CCF-A类顶级安全会议}

\begin{itemize}
  \item 提出Fuzzer-in-the-Middle网络协议模糊测试架构，首次系统地测试Linux内核网络栈远程攻击面，发现\textbf{15+}个远程可触发内核漏洞。
  \item 负责核心架构设计与实现、LibAFL定制、漏洞复现、实验评估及论文主要撰写。
\end{itemize}

\entrydivider

\pub{CARDSHARK: Understanding and Stabilizing Linux Kernel Concurrency Bugs Against the Odds}
{\textbf{Tianshuo Han}, Xiaorui Gong, and Jian Liu}
{33rd USENIX Security, 2024, CCF-A类顶级安全会议}

\begin{itemize}
  \item 分析Linux内核竞争漏洞触发不稳定的根因，提出了一种通用的竞争漏洞利用技术，显著提升漏洞利用成功率和可复现性。
  \item 负责竞争漏洞机理分析、利用技术研发、性能评估与论文主要撰写。
\end{itemize}

\entrydivider

\pub{Venom: Kernel Fuzzing of Distributed File Systems via Intra-Cluster Protocol Poisoning}
{Zong Cao， \textbf{Tianshuo Han}（共同一作），等}
{ACM CCS, 2026，CCF-A类顶级安全会议（在投）}

\begin{itemize}
  \item 提出Mutation-in-the-Middle及静态分析引导变异算法，在内核分布式文件系统发现\textbf{30+}个内核漏洞。
  \item 负责变异算法设计与实现、实验设计、漏洞复现及结果分析。
\end{itemize}

\entrydivider

\pub{Wire to Ring Zero: The Remote Kernel Attack Surface That Years of Fuzzing Never Reached}
{\textbf{Tianshuo Han}，Zong Cao}
{Black Hat USA，2026，顶级安全工业会议（在投）}

\begin{itemize}
  \item 从工业攻防视角梳理远程内核攻击面，提出面向Linux/Windows内核网络栈实现的中间人模糊测试方法。
\end{itemize}

\entrydivider

\pub{Reviving Discarded Vulnerabilities: Exploiting Previously Unexploitable Linux Kernel Bugs Through Control Metadata Fields}
{Hao Zhang, Jie Lu, Shaomin Chen, \textbf{Tianshuo Han}, Bolun Zhang, Jian Liu, Xiaorui Gong}
{32nd CCS, 2025, CCF-A类顶级安全会议}

\begin{itemize}
  \item 提出基于内核结构体关键元数据字段操控的利用技术，增强弱漏洞原语场景下的利用能力，扩展可利用漏洞范围。
  \item 参与核心思路制定、评估实验设计及审稿意见回复。
\end{itemize}

\entrydivider

\noindent\textbf{Whole-genome sequence and comparative analysis of \emph{Trichoderma asperellum} ND-1}\par
\noindent{\small Fengzhen Zheng, \textbf{Tianshuo Han}, et al. \enspace|\enspace Catalysts 12(4), 2022, IF=4.16}\par

\begin{itemize}
  \item 负责基因组数据分析与科学数据可视化，辅助揭示该菌株木质纤维素降解酶系统。
\end{itemize}

\newpage

% ── 核心项目 ──
\section{项目实践}

\resumeentry{NetPanic：面向Linux内核网络栈的FitM模糊测试框架}{2023.11 -- 至今}

\begin{itemize}
  \item 设计区别于Syzkaller的在途中间人模糊测试架构，面向NFS/CIFS/NFSD/ksmbd等网络路径提升远程攻击面覆盖。
  \item 基于LibAFL深度定制变异、执行、反馈、调度等核心组件，独立完成\textbf{40K+}行Rust/C实现。
  \item 支撑IEEE S\&P 2026论文与多项远程可触发Linux内核CVE发现。
\end{itemize}

\entrydivider

\resumeentry{Venom：Mutation-in-the-Middle变异算法}{2026.03 -- 至今}

\begin{itemize}
  \item 设计静态分析引导的高效变异算法，针对Fuzzer-in-the-Middle架构中的网络协议状态、消息依赖和污染传播进行优化。
  \item 基于LibAFL与LLVM IR实现，代码规模\textbf{50K+}行（Rust/C/Python），支撑分布式文件系统漏洞发现。
\end{itemize}

\entrydivider

\resumeentry{自动化漏洞审计与复现Agent}{2026.03 -- 2026.04}

\begin{itemize}
  \item 基于ClaudeCode/Codex SDK构建端到端全自动安全审计Agent，覆盖漏洞发现、环境搭建、复现、PoC生成和报告撰写。
  \item 在HTTPD、QEMU、libexif等开源项目中发现\textbf{10+}个漏洞并获得CVE。
\end{itemize}

\entrydivider
\resumeentry{Linux内核Rootkit与隐蔽机制研究}{2023.12 -- 2024.05}

\begin{itemize}
  \item 独立实现自身隐藏、端口隐藏、权限提升、文件/进程隐藏、开机自启等核心机制，用于内核攻防技术验证。
  \item 强调轻量化、高隐蔽性与跨版本兼容，支持Linux 2.6至6.3系列内核并规避常见Rootkit检测策略。
\end{itemize}

\entrydivider

\resumeentry{类Linux内核开发}{2020.12 -- 2021.05}

\begin{itemize}
  \item 从零实现可运行的类Linux内核，覆盖引导、内存管理、进程调度、系统调用和文件系统等核心子系统；代码规模约\textbf{20K}行（C/x86汇编）。
\end{itemize}

% ── 安全贡献 ──
\section{安全贡献}

\resumeentry{Linux内核远程可触发漏洞}{5项}

\begin{itemize}
  \item \textbf{CVE-2025-38089}：sunrpc空指针解引用，可远程触发内核崩溃
  \item \textbf{CVE-2025-38051}：CIFS客户端\texttt{cifs\_fill\_dirent}释放后重用
  \item \textbf{CVE-2025-40210}：NFSD远程COMPOUND请求内存分配失败拒绝服务
  \item \textbf{CVE-2025-40212}：NFSD伪根文件句柄引用计数泄漏，导致UAF与拒绝服务
  \item \textbf{CVE-2026-23220}：ksmbd SMB2请求处理无限循环远程拒绝服务
\end{itemize}

\vspace{2pt}

\resumeentry{开源项目漏洞}{5项}

\begin{itemize}
  \item \textbf{CVE-2026-28780}、\textbf{CVE-2026-34032}：暂未公开披露
  \item \textbf{CVE-2026-40312}：ImageMagick MSL解码器XML分组解析堆越界写入（off-by-one）
  \item \textbf{CVE-2026-40385}：libexif Nikon MakerNote 32位无符号整数溢出，导致崩溃或信息泄露
  \item \textbf{CVE-2026-40386}：libexif Fuji/Olympus MakerNote尺寸校验整数下溢，导致崩溃或信息泄露
\end{itemize}

% ── 技术能力 ──
\section{技术能力与其他}

\begin{itemize}
  \item \textbf{内核安全}：Linux内核漏洞挖掘与利用、远程攻击面分析、Rootkit机制、二进制逆向与利用。
  \item \textbf{安全能力}：熟练掌握网络、系统和密码等网络安全基础；熟悉LibAFL、AFL++、Syzkaller、LLVM、CodeQL等模糊测试和静态分析框架；熟练使用GDB、IDA等程序调试和逆向工具；熟练使用各种渗透测试工具。
  \item \textbf{编程语言}：精通\textbf{C}和\textbf{Rust}；熟练掌握Python、Shell、x86汇编、\LaTeX{}，可以使用各种主流编程语言。
  \item \textbf{AI与Harness}：熟练掌握Agentic编程工作流，可以快速完成开发任务和安全测试任务。
  \item \textbf{教学经历}：曾任中国科学院大学《恶意软件发现与分析》课程助教；
  \item \textbf{外语能力}：英语读写能力与母语者接近，可熟练阅读、撰写英文论文和技术文档。
\end{itemize}

\end{document}
